% Methodology / Technical Route
\section{Methodology and Technical Route}
\label{sec:method}

This section specifies the proposed Long-Range Video Consistency Composite (LR-VCC) metric,
its design rationale, and the iterative protocol under which the current production version
was derived. The metric is constructed directly from the four failure modes characterised
in Section~\ref{sec:problem}; each sub-metric closes a specific gap, and the composition
layer downweights any sub-metric whose calibration assumptions are violated on the video
under evaluation.

\subsection{Architecture}

\textbf{[Figure 4: LR-VCC architecture --- super-resolved video as input; five sub-metrics
computed in parallel, each emitting a $(\text{score}, \text{reliability})$ pair in
$[0, 1] \times [0, 1]$; a reliability-weighted softmax-log-mean composition produces the
final LR-VCC score in $[0, 1]$.]}

Five sub-metrics, indexed by $s \in \{A, T, I, D, E\}$, are computed in parallel from a
super-resolved video $V$:

\begin{itemize}
    \item \textbf{Sub-metric A --- Appearance stability.} Per-frame CLIP-IQA
          \cite{wang_clipiqa_2023} is computed over the video. The score is the mean
          per-frame quality minus a small multiple of its standard deviation across frames,
          rewarding outputs that are not only high-quality on average but consistently so
          across minutes of footage. The mean term rewards perceptual quality; the standard
          deviation term penalises temporal instability of perceptual quality.

    \item \textbf{Sub-metric T --- Temporal stability.} Optical-flow warping error (tOF
          \cite{chu_tecogan_2020}) is computed at multiple temporal lags
          $k \in \{1, 5, 10, 30, 60, 120\}$ using RAFT-based optical flow with
          forward-backward consistency masking. The aggregate is a weighted mean over $k$,
          with the choice of weighting (logarithmic, uniform, or square-root) configurable;
          uniform weighting is adopted as the production setting because the empirical
          temporal-scale crossover characterised in Section~\ref{sec:problem:tof} shows
          that no single $k$ is sufficient. The score is one minus the (normalised)
          weighted mean warping error.

    \item \textbf{Sub-metric I --- Identity preservation.} The slow-fast Human\_Identity
          adapter (RetinaFace + ArcFace face embedding stability across $2$-second clips
          with both within-clip and cross-clip first-frame components) provides the score,
          with reliability gated by the face-detection rate (a video with few detected
          faces does not carry an identity signal) and by a close-up content indicator
          (the same trigger identified in Section~\ref{sec:problem:regime}) that
          downweights the sub-metric on videos in the regime where the underlying anomaly
          classifiers exhibit content-dependent calibration shifts.

    \item \textbf{Sub-metric D --- Colour stability.} Per-frame Lab-channel histograms are
          compared at L1 distance over multi-scale frame pairs $k \in \{60, 120\}$. The
          score is an exponential decay in the mean pair-distance, with the decay constant
          calibrated empirically so that clean super-resolved outputs score in the middle
          of the dynamic range. Reliability is gated by histogram entropy --- a video
          whose histogram is too flat carries no discriminative colour information.

    \item \textbf{Sub-metric E --- Colour trajectory slope.} Per-frame Lab-channel means
          are fit by linear regression as a function of frame index. The score is an
          exponential decay in the maximum absolute slope across channels. Reliability is
          the maximum coefficient of determination ($R^2$) across the three channels, gated
          by an $R^2$ floor. A drifting video has high $R^2$ (the slope is real, so the
          metric speaks confidently); a flickering video has $R^2$ near zero (sinusoidal
          residuals dominate, so the metric correctly abstains); a clean video has
          near-zero slope (the score is near one regardless of reliability).
\end{itemize}

Each sub-metric outputs a score $s_i \in [0, 1]$ and a reliability $r_i \in [0, 1]$. The
composition is a reliability-weighted log-mean:
\begin{align}
    w_i &= \frac{\exp(r_i / \tau)}{\sum_j \exp(r_j / \tau)}, \\
    \text{LR-VCC}(V) &= \exp\!\left(\sum_{i} w_i \log(s_i + \varepsilon)\right),
\end{align}
with temperature $\tau = 0.2$ (a sharp softmax: the most reliable sub-metric dominates) and
$\varepsilon = 10^{-6}$ for numerical stability. A video for which all five reliabilities
fall below a low-confidence threshold is flagged in the output and excluded from
per-method aggregates by default.

The log-mean composition preserves the no-compensation property across sub-metrics: a
sub-metric scoring near zero pulls the composite down unless the other sub-metrics carry
overwhelming weight, which in turn requires the failing sub-metric's reliability to be near
zero. A weighted arithmetic mean would not have this property: a single sub-metric failure
could be diluted away by averaging.

\subsection{Reliability Gating as the Core Design Pattern}

The central design hypothesis is that every existing learned-representation video metric is
reliable on some content and unreliable on other content; rather than discard failed
metrics or attempt to train new bias-free metrics, the composition uses existing metrics
honestly and controls their influence per video through reliability scores derived from
the regime indicators that predict their failure mechanisms.

Each per-sub-metric reliability is computed from regime indicators using smooth sigmoid
penalties (sharpness fixed at $10$, producing smooth transitions rather than sharp threshold
cliffs). The thresholds were selected from the regime characterisations in
Section~\ref{sec:problem}; they are not tuned to the validation set used in
Section~\ref{sec:method:iteration}.

A sharp softmax over reliabilities ensures that when one sub-metric is substantially more
trustworthy than the others on a given video, that sub-metric clearly dominates the
composition. This is the mechanism that makes the composite resistant to the
content-dependent ranking flips of Section~\ref{sec:problem:regime}: on a video in the
close-up regime, where the Identity-pathway anomaly classifier is known to misbehave, the
Identity sub-metric's reliability is downweighted by the close-up gate, allowing the other
four sub-metrics to determine the composite outcome.

\subsection{Validation Strategy}

The composite is validated at three layers, as introduced in
Section~\ref{sec:obj:validate}.

\textbf{Layer~1 --- perceptual agreement on the real-SR test set.} The composite is
evaluated on the five-video test set comprising both super-resolution baselines. The pass
criterion is that the detail-preserving baseline wins on the per-method mean and on every
per-video comparison. The current production version meets this criterion: the per-video
comparison is consistent across all test videos including the regime-shift case, with a
mean difference between the two methods of approximately $+0.06$ in favour of the
detail-preserving baseline.

\textbf{Layer~2 --- flip-resistance on the regime-shift case.} The composite must not flip
on the close-up video where Human\_Anatomy and Human\_Identity both flip. The criterion is
met by the design: sub-metric I's reliability gate engages on close-up content, sub-metric
T uses warping error directly (not LPIPS-based tLP, which exhibits the smoother-output bias
of its underlying perceptual representation), and sub-metric A is built on CLIP-IQA, which
does not depend on a pristine-HR-trained representation.

\textbf{Layer~3 --- severity-response on the synthetic artefact set.} For each of the four
artefact families, the composite is expected to respond monotonically to severity on at
least one base video. The current production version catches the majority of (artefact
family, base video) cells monotonically. The remaining cells have characterised failure
modes:
\begin{itemize}
    \item On the single-face base video, identity degradation triggers the
          identity-collapse pathology of Section~\ref{sec:problem:battery} Finding~4. The
          face-detection rate remains high, so the existing face-rate reliability gate does
          not engage; the pathology lies in the slow-fast pooling itself. A reliability
          gate that operates on per-face embedding variance is identified as the natural
          fix (Section~\ref{sec:outcomes}).
    \item Flicker is correctly detected by sub-metric T at small $k$, but the composite is
          dominated by the colour and identity sub-metrics whose reliabilities are near
          one and whose scores do not respond to high-frequency luminance modulation; the
          composite is therefore approximately flat across flicker severities. A
          fast-varying brightness sub-metric is the natural addition.
    \item One of the two base videos for colour drift exhibits a pre-existing slow colour
          trajectory in the underlying super-resolved output. The slope-style sub-metric
          correctly assigns high reliability to that pre-existing trajectory, leaving
          little additional dynamic range to register the injected drift. The case is
          informative rather than spurious.
\end{itemize}

The full verdict matrix and per-sub-metric severity-response curves are presented in
Figures~5--7.

\subsection{Iterative Design Protocol}
\label{sec:method:iteration}

The current production version of the composite is the result of an iterative protocol in
which each version added exactly one element after a specific failure mode was characterised
on the synthetic test set. The history is summarised in
Table~\ref{tab:method:iterations}.

\begin{table}[h]
\centering
\caption{Iterative design history of the composite. Each row introduces one element after a
specific failure mode was characterised on the synthetic artefact set; the next row
addresses the next failure surfaced by the updated composite.}
\label{tab:method:iterations}
\small
\begin{tabular}{p{2.8cm}p{5.5cm}p{6.0cm}}
\hline
\textbf{Version} & \textbf{New element} & \textbf{Failure mode closed} \\
\hline
v1 & Three sub-metrics: Appearance, Temporal (logarithmic-weighted), Identity & Initial design responsive to the regime-shift case and the temporal-scale crossover \\
v2 & Add Colour Stability sub-metric (Lab histogram L1 over multi-scale pairs) & Provides a long-range colour lever distinct from optical-flow-based metrics \\
v2-uniform & Change temporal weighting to uniform & The high-frequency component at small $k$ is no longer drowned out at the sub-metric T level \\
v3 & Recalibrate Colour Stability decay constant & The colour sub-metric's absolute score range no longer dominates the composite arithmetic; chunk-boundary signal flows through to the composite \\
v3-slope & Add Colour Trajectory Slope sub-metric with goodness-of-fit-gated reliability & Slow colour drift is detected (Section~\ref{sec:problem:battery} Finding~1) on at least one base video \\
\hline
\end{tabular}
\end{table}

The protocol itself --- characterise the failure mode, add one element, verify on the same
test set without regression on previous failure modes --- is reusable beyond the present
metric. Any composite no-reference metric whose failure modes can be parameterised as
controlled synthetic artefacts admits an analogous iteration, with the goodness-of-fit-gated
reliability pattern of sub-metric E reusable for any ``slow signal versus noisy baseline''
detection problem in video evaluation.

\subsection{Reproducibility and Decomposability}

The metric is fully decomposable. Each sub-metric is computed independently, and the result
is cached per video. Hyperparameter studies and re-tuning of the per-sub-metric decay
constants and weighting choices do not require re-scanning videos; the cached per-video
values are re-aggregated through the composition layer in seconds. This property is
essential to the iterative protocol of
Section~\ref{sec:method:iteration}: each version can be tested against the same cache
without paying the cost of recomputing the sub-metrics.

The output of the composite on a given video is the LR-VCC score, the per-sub-metric
$(\text{score}, \text{reliability})$ pairs, the softmax-derived sub-metric weights, the
low-confidence flag, and a set of per-sub-metric diagnostic indicators (closeup indicator,
face-detection rate, mean mask coverage at long temporal lags, regression $R^2$, histogram
entropy). The per-video per-sub-metric matrix is always reported alongside the headline
composite, so a reviewer can see where the composite is driven from on each video. There is
no black-box aggregation.
